\documentclass[11pt]{article}
\usepackage[T1]{fontenc}
\usepackage{lmodern}
\usepackage[margin=1.05in]{geometry}
\usepackage{amsmath,amssymb,amsthm,mathtools}
\usepackage{microtype,booktabs,array}
\usepackage{listings}
\usepackage{xcolor}
\usepackage[colorlinks=true,linkcolor=blue!45!black,citecolor=blue!45!black,urlcolor=blue!45!black]{hyperref}
\usepackage[nameinlink,capitalise,noabbrev]{cleveref}
\numberwithin{equation}{section}
\newtheorem{theorem}{Theorem}[section]
\newtheorem{proposition}[theorem]{Proposition}
\newtheorem{lemma}[theorem]{Lemma}
\newtheorem{corollary}[theorem]{Corollary}
\newtheorem{conjecture}[theorem]{Conjecture}
\theoremstyle{definition}
\newtheorem{definition}[theorem]{Definition}
\theoremstyle{remark}
\newtheorem{remark}[theorem]{Remark}
\newcommand{\Z}{\mathbb Z}
\newcommand{\F}{\mathbb F}
\newcommand{\CW}{\mathrm{CW}}
\newcommand{\ICW}{\mathrm{ICW}}
\DeclareMathOperator{\supp}{supp}
\DeclareMathOperator{\Tr}{T}
\DeclareMathOperator{\Nm}{N}
\DeclareMathOperator{\PAF}{PAF}
\lstdefinelanguage{lean}{morekeywords={def,theorem,instance,by,fun,if,then,else,let,example},sensitive=true,
  morecomment=[l]{--},morecomment=[s]{/-}{-/},morestring=[b]"}
\lstset{language=lean,basicstyle=\ttfamily\footnotesize,columns=fullflexible,keepspaces=true,
  keywordstyle=\bfseries,commentstyle=\itshape\color{black!60},breaklines=true,
  literate={∀}{{$\forall$}}1 {∃}{{$\exists$}}1 {∧}{{$\wedge$}}1 {∨}{{$\vee$}}1 {→}{{$\to$}}1
           {≤}{{$\le$}}1 {≠}{{$\ne$}}1 {∣}{{$\mid$}}1 {∈}{{$\in$}}1 {¬}{{$\neg$}}1 {Σ}{{$\Sigma$}}1
           {·}{{$\cdot$}}1 {⟨}{{$\langle$}}1 {⟩}{{$\rangle$}}1}
\title{New proper circulant weighing matrices\\ from a twin of the Singer construction}
\author{Claude (AI model by Anthropic)\\ \small run and submitted by Aleksandr Bersenev}
\date{October 8, 2026\\[3pt]\footnotesize Produced autonomously by Claude (Opus 5.5) agents in Claude Code; run and submitted by Aleksandr Bersenev.\\ \footnotesize Not an official Anthropic publication. Not peer reviewed.}
\hypersetup{pdftitle={New proper circulant weighing matrices from a twin of the Singer construction},pdfauthor={Claude (AI model by Anthropic); submitted by Aleksandr Bersenev}}
\begin{document}
\maketitle
\begin{abstract}
A circulant weighing matrix $\CW(n,k)$ is a circulant $n\times n$ matrix $W$ with entries in
$\{0,\pm1\}$ and $WW^{\mathsf T}=kI$; it is proper if it is not the lift of a smaller one.
We construct proper
$\CW(146,8^2)$, $\CW(266,11^2)$, $\CW(366,13^2)$, $\CW(504,12^2)$, $\CW(614,17^2)$ and
$\CW(762,19^2)$. To our knowledge, as of October 8, 2026, the existence of a proper
circulant weighing matrix with each of these parameters was open: all six are listed as open
in Gordon's La Jolla Circulant Weighing Matrix Repository and none appears in Gordon's 2026
table of known proper $\CW(n,k)$ with $n<1000$, $k\le 19^2$. For $\CW(504,144)$ no circulant
weighing matrix of any kind was known.
All six first rows are explicit and are verified in Lean~4 (core library only) by kernel
evaluation.
Four of them come from a new ``twin'' of the classical Singer-complement construction.
For $q$ odd, the standard $\CW(q^2+q+1,q^2)$ on the points of $\mathrm{PG}(2,q)$ has
a twin, obtained by composing with the homology of ratio $-2$ (an elation when $p=3$), that has
the same support.
Whenever the twin is again a circulant weighing matrix, a doubling lemma produces a proper
$\CW(2(q^2+q+1),q^2)$.
We conjecture that this holds for every odd prime power $q$. It has been checked by computer for
all odd prime powers $q\le251$, but we have no proof.
The other two matrices come from a multiplier-restricted search ($n=146$) and from a product
over non-coprime orders ($n=504$).
\end{abstract}
\tableofcontents

\section{Introduction}\label{sec:intro}

A \emph{circulant weighing matrix} $\CW(n,k)$ is an $n\times n$ circulant matrix with entries in
$\{0,1,-1\}$ such that $WW^{\mathsf T}=kI_n$. It is determined by its first row
$a=(a_0,\dots,a_{n-1})$, and the condition says exactly that $a$ is a \emph{perfect ternary
sequence}: its periodic autocorrelation
\[
  \PAF_a(d)=\sum_{i=0}^{n-1}a_i\,a_{(i+d)\bmod n}
\]
equals $k$ for $d=0$ and $0$ for $0<d<n$. Necessarily $k=s^2$ is a square. If $W$ is a
$\CW(m,k)$ and $m\mid n$, then spreading its first row out (the group-ring map $X\mapsto
X^{n/m}$) gives a $\CW(n,k)$. Matrices that do not arise in this way are called
\emph{proper}: the support of the first row is not contained in a coset of a proper subgroup of
$\Z_n$.

Strassler's 1997 table of existence results for $n\le200$, $k\le100$ was updated by Tan and
reduced to 27 open cases by Arasu, Gordon and Zhang~\cite{AGZ}. Tan~\cite{Tan26} recently
settled eight more cases (all nonexistence results).
Gordon maintains the La Jolla Circulant Weighing Matrix Repository~\cite{GordonRepo}. It records
the existence status of proper $\CW(n,s^2)$ for all $n\le1000$, $s\le19$, together with explicit
matrices.
Table~7 of~\cite{G26}, reproduced on the repository web page (last updated 2026-09-03), lists
all \emph{known} proper $\CW(n,k)$ with $n<1000$ and $k\le19^2$. The web page remarks that for
larger $k$ almost all entries come from the product or relative-difference-set (RDS)
constructions, ``so presumably many others remain to be found''.
For odd prime-power weight $k$, Leung and Schmidt~\cite{LS11} proved that only finitely many
proper $\CW(n,k)$ exist (for $k=4$ there are infinitely many, so the parity condition is needed).
Leung and Ma~\cite{Ma11} asked (open problem~7 in the slides of a 2011 talk): apart from the classical construction of proper
$\CW((p^2+p+1)w/2,p^2)$ ($w\mid p-1$, $w\equiv2\pmod 4$), are there other (infinite)
families of proper $\CW(v,p^2)$?

A further piece of history is relevant here. Theorem~6.6 of~\cite{AGZ} asserted the existence of a
proper $\CW((q^3-1)/n,q^2)$ for every divisor $n>1$ of $q-1$ when $q$ is odd. Its proof used an
RDS-to-CW theorem (Theorem~6.4 there) which, as Gordon later pointed out~\cite[Theorem~12]{G26},
needs both indices $m$ and $n$ of the $(m,2n,k,\lambda)$-RDS to be odd.
Gordon's corrected version (Theorem~13 of~\cite{G26}) only yields orders $(q^2+q+1)n'$ with $n'$
odd. Accordingly, the 2025/26 table no longer lists $366$ $(13^2)$, $614$ $(17^2)$ or $762$
$(19^2)$, all of which were listed in Table~11 of~\cite{AGZ}. The case $n=(q-1)/2$ of the
withdrawn claim is precisely the family $\CW(2(q^2+q+1),q^2)$ studied here.

\paragraph{Results.} Table~\ref{tab:results} lists the six parameter pairs settled here.

\begin{table}[h]
\centering\small
\setlength{\tabcolsep}{4pt}
\begin{tabular}{@{}lllll@{}}
\toprule
$(n,k)$ & repository~\cite{GordonRepo} & Table~7 of~\cite{G26} & non-proper CW known? & construction\\
\midrule
$(146,8^2)$  & Open & no & yes (lift of $\CW(73,8^2)$) & multiplier search, \S\ref{sec:146}\\
$(266,11^2)$ & Open & no & yes (lift of $\CW(133,11^2)$) & twin, $q=11$, \S\ref{sec:twin}\\
$(366,13^2)$ & Open & no (dropped) & yes (lift of $\CW(183,13^2)$) & twin, $q=13$\\
$(504,12^2)$ & Open & no & no & product, \S\ref{sec:504}\\
$(614,17^2)$ & Open & no (dropped) & yes (lift of $\CW(307,17^2)$) & twin, $q=17$\\
$(762,19^2)$ & Open & no (dropped) & yes (lift of $\CW(381,19^2)$) & twin, $q=19$\\
\bottomrule
\end{tabular}
\caption{The six parameter pairs. ``Open'' is the status in the repository data file (v1.3, commit
5b3d729, 2026-04-24); the second parameter there is $s=\sqrt k$. ``Dropped'' means the entry
appeared in~\cite[Table~11]{AGZ} but not in~\cite[Table~7]{G26}.}
\label{tab:results}
\end{table}

\begin{theorem}[main theorem; Lean-certified]\label{thm:main}
There exist proper circulant weighing matrices $\CW(146,64)$, $\CW(266,121)$, $\CW(366,169)$,
$\CW(504,144)$, $\CW(614,289)$ and $\CW(762,361)$.
\end{theorem}

\begin{conjecture}[twin conjecture; Conjecture~\ref{conj:twin} below]
For every odd prime power $q$ the twin $W_{1/2}$ of the classical $\CW(q^2+q+1,q^2)$ is again a
circulant weighing matrix. Consequently a proper $\CW(2(q^2+q+1),q^2)$ exists for every odd
prime power $q$.
\end{conjecture}

The conjecture has been verified for every odd prime power $q\le251$ (62 values; see
Remark~\ref{rem:evidence}). For
$q=3,5,7,9$ it recovers the orders $26,62,114,182$, which~\cite{G26} lists as sporadic (boxed)
entries. If true, it would answer the question of Leung and Ma with a new infinite family. We stress that we do
\emph{not} have a proof.

\paragraph{Novelty and caveats.} We compared with~\cite{GordonRepo,G26,AGZ,Tan26,LS11,SS13} and
with the perfect-ternary-sequence literature that we could access (Ipatov's construction as
described in~\cite{HJ83}, and~\cite{HJ83} itself). We could \emph{not} check
Krengel~\cite{Kr04}, which (according to its abstract) constructs almost-perfect ternary sequences
of length $(p^n-1)/r$ and odd-perfect ternary sequences of half that length from $m$-sequences, in
particular odd-perfect sequences of length $2(p^n-1)/(p^{n/3}-1)=2(q^2+q+1)$. Odd-perfect sequences
are not perfect, so these are different objects from our $\CW(2(q^2+q+1),q^2)$; but a perfect
sequence of odd length $m$ becomes odd-perfect after the sign change $i\mapsto(-1)^i$, so the twin
$W_{1/2}$ itself might appear there in another guise. Our novelty claims are therefore ``to our
knowledge, as of October 8, 2026''.

\section{Preliminaries}\label{sec:prelim}

We identify $\Z_n$ with the cyclic group $\langle X\rangle$, $X^n=1$, and a sequence $a$ with
$A=\sum_i a_iX^i\in\Z[\Z_n]$. For $A=\sum_g a_g g$ in a group ring $\Z[G]$ of a finite abelian
group $G$ we write $A^{(-1)}=\sum_g a_g g^{-1}$ and $\supp A=\{g:a_g\ne0\}$.
Then $a$ is the first row of a $\CW(n,k)$ if and only if $a_i\in\{0,\pm1\}$ and
$AA^{(-1)}=k$. We say $A$ is a $\CW$ in $G$ if this holds with $G$ in place of $\Z_n$.

\begin{definition}
$A\in\Z[G]$ is \emph{proper} if $\supp A$ is not contained in a coset $g+L$ of a proper
subgroup $L<G$. For $G=\Z_n$ the subgroups are $e\Z_n$ ($e\mid n$), so $A$ is proper if and
only if for every divisor $e>1$ of $n$ the support meets at least two residue classes modulo $e$.
\end{definition}

\begin{lemma}[Fourier criterion]\label{lem:fourier}
For $A\in\Z[G]$ we have $AA^{(-1)}=k$ if and only if $|\chi(A)|^2=k$ for every character $\chi$
of $G$.
\end{lemma}
This is standard: $\chi(AA^{(-1)})=|\chi(A)|^2$, and an element of $\mathbb C[G]$ is determined
by its character values.

\paragraph{Field notation.} Let $q$ be an odd prime power, $F=\F_q$, $K=\F_{q^3}$, and let
$\Tr,\Nm\colon K\to F$ be the trace and norm. Let $\eta$ be the quadratic character of $F$
(with $\eta(0)=0$) and $\eta_K=\eta\circ \Nm$; this is the quadratic character of $K$, because
$\Nm\colon K^*\to F^*$ is onto. Put $m=q^2+q+1$ and $G=K^*/F^*$, a cyclic group of order $m$.
Its elements are the points $[y]$ of $\mathrm{PG}(2,q)$, where $K$ is viewed as $F^3$.
Fix a primitive element $\alpha$ of $K$ and identify $i\in\Z_m$ with $[\alpha^i]$; under this
identification multiplication in $K^*$ becomes the Singer cycle.
Let $L=\{[y]:\Tr(y)=0\}$, a line of $q+1$ points (the classical Singer difference set), and let
$S=G\setminus L$, so $|S|=q^2$. Every point of $S$ has exactly one representative in the affine
plane $P=\{y\in K:\Tr(y)=1\}$.

\section{Three elementary lemmas}\label{sec:lemmas}

\begin{lemma}[doubling]\label{lem:doubling}
Let $m$ be odd. Identify $\Z_{2m}\cong\Z_m\times\Z_2$ by $x\mapsto(x\bmod m,\,x\bmod 2)$, and
let $z$ be the element of order $2$. For $U,V\in\Z[\Z_m]$ put $A=U+zV$. Then $AA^{(-1)}=k$ if
and only if $(U+V)(U+V)^{(-1)}=k$ and $(U-V)(U-V)^{(-1)}=k$.
\end{lemma}
\begin{proof}
Every character of $\Z_m\times\Z_2$ is $\chi\otimes\epsilon$ with $\epsilon(z)=\pm1$, and
$(\chi\otimes\epsilon)(A)=\chi(U\pm V)$. Apply Lemma~\ref{lem:fourier} in $\Z_{2m}$ and in
$\Z_m$.
\end{proof}

\begin{corollary}\label{cor:samesupport}
Let $m$ be odd and let $W,W'$ be first rows of two $\CW(m,k)$ with the same support. Define
$A\in\{0,\pm1\}^{2m}$ by $A_x=\tfrac12(W+W')_{x\bmod m}$ for even $x$ and
$A_x=\tfrac12(W-W')_{x\bmod m}$ for odd $x$. Then $A$ is a $\CW(2m,k)$.
If in addition $W$ is proper and $W'\ne\pm W$, then $A$ is proper.
\end{corollary}
\begin{proof}
Since $W'_i=\pm W_i$ on the common support, $U=\frac12(W+W')$ and $V=\frac12(W-W')$ have
entries in $\{0,\pm1\}$ and disjoint supports. Under the identification of
Lemma~\ref{lem:doubling}, $A=U+zV$, and $U\pm V\in\{W,W'\}$, so $A$ is a $\CW(2m,k)$.

For properness, let $e\mid 2m$ with $e>1$ and suppose $\supp A\subseteq c+e\Z_{2m}$.
If $e$ has an odd prime factor $p$, then $p\mid m$ and every $x\in\supp A$ satisfies
$x\equiv c\pmod p$. Hence $x\bmod m$ lies in one residue class mod $p$, so $\supp W$ lies in a
coset of $p\Z_m$, contradicting properness of $W$. Otherwise $e=2$ (as $4\nmid 2m$), so
$\supp A$ lies in one parity class, i.e.\ $U=0$ or $V=0$, i.e.\ $W'=\pm W$.
\end{proof}

\begin{lemma}[product over non-coprime orders]\label{lem:product}
Let $H,K'$ be subgroups of a finite abelian group $G$, and let $A\in\Z[H]$, $B\in\Z[K']$ be
$\CW$s in $H$ and $K'$ of weights $k_1,k_2$. Assume
\[
(\supp A-\supp A)\cap(\supp B-\supp B)=\{0\}.
\]
Then $AB$ has coefficients in $\{0,\pm1\}$ and is a $\CW$ of weight $k_1k_2$ in $G$. If moreover
$A$ is proper in $H$, $B$ is proper in $K'$ and $H+K'=G$, then $AB$ is proper in $G$.
\end{lemma}
\begin{proof}
If $h+b=h'+b'$ with $h,h'\in\supp A$ and $b,b'\in\supp B$, then $h-h'=b'-b$ lies in both
difference sets, so $h=h'$ and $b=b'$. Hence the $|\supp A|\cdot|\supp B|$ products are
distinct, and every coefficient of $AB$ is $0$ or $\pm1$. Also
$(AB)(AB)^{(-1)}=AA^{(-1)}BB^{(-1)}=k_1k_2$.

Now suppose $\supp(AB)\subseteq c+L$ for some $L<G$. Fixing $b\in\supp B$ gives
$\supp A\subseteq c-b+L$, hence $\supp A-\supp A\subseteq L\cap H$. Since $A$ is proper in $H$,
this forces $L\cap H=H$, i.e.\ $H\le L$. Likewise $K'\le L$, so $L\supseteq H+K'=G$.
\end{proof}

For coprime orders, Lemma~\ref{lem:product} is the Kronecker product construction of Arasu and
Seberry~\cite{AS98} (where $H\cap K'=0$). The version for non-coprime orders is elementary and
may well be folklore. We do not claim it as new; we only note that it was not applied to
$n=504$ in the tables.

\section{The twin of the Singer-complement construction}\label{sec:twin}

\begin{definition}\label{def:Wc}
For $c\in F$ define $W_c\colon G\to\{0,\pm1\}$ by
\[
  W_c([y])=\eta\bigl(\Nm(c\,\Tr(y)-y)\cdot \Tr(y)\bigr).
\]
This is well defined: replacing $y$ by $\lambda y$ ($\lambda\in F^*$) multiplies the argument of
$\eta$ by $\lambda^4$. As usual, $W_c$ also denotes $\sum_{g}W_c(g)\,g\in\Z[G]$ and, via
$i\mapsto[\alpha^i]$, a sequence of length $m$.
\end{definition}

\begin{lemma}\label{lem:support}
$\supp W_0=\supp W_{1/2}=S$.
\end{lemma}
\begin{proof}
Both values vanish when $\Tr(y)=0$. Let $\Tr(y)\neq0$. Then $\Nm(-y)\neq0$. If
$\tfrac12\Tr(y)-y=0$, then $y\in F^*$, so $\Tr(y)=3y$ and $\tfrac12\Tr(y)-y=\tfrac12y\ne0$, a
contradiction.
\end{proof}

\begin{proposition}[classical]\label{prop:classical}
$W_0$ is a proper $\CW(m,q^2)$.
\end{proposition}
\begin{proof}
For $y\in P$ we have $W_0([y])=\eta(\Nm(-y))$. Let $\delta\in K^*\setminus F^*$. Using the
homogeneous formula for $W_0([\delta y])$ we get
\[
\sum_{y\in P}W_0([y])W_0([\delta y])
 =\sum_{y\in P}\eta(\Nm(-y))\,\eta\bigl(\Nm(-\delta y)\Tr(\delta y)\bigr)
 =\eta(\Nm\delta)\sum_{y\in P}\eta(\Tr(\delta y)).
\]
The map $y\mapsto \Tr(\delta y)$ is affine on the affine plane $P$. It is not constant: otherwise
$\Tr(\delta z)=0$ for all $z\in\ker \Tr$, so $\Tr(\delta\,\cdot)=\lambda \Tr$ for some
$\lambda\in F$, and nondegeneracy of the trace form gives $\delta=\lambda\in F$. A
nonconstant affine function on $\F_q^2$ takes every value exactly $q$ times, and
$\sum_{t\in F}\eta(t)=0$. So all nontrivial autocorrelations vanish, and the trivial one equals
$|S|=q^2$. Finally, a coset of a proper subgroup of $\Z_m$ has at most $m/3<q^2$ elements, so
$W_0$ is proper.
\end{proof}

This is a standard construction of a $\CW(q^2+q+1,q^2)$; cf.~\cite[Theorem~4]{AS98} and the
RDS constructions in~\cite{G26}. Our object of interest is $W_{1/2}$.

\begin{remark}[geometry]\label{rem:geometry}
Let $\psi(y)=\Tr(y)-2y$. This map is $F$-linear and invertible, with
$\psi^{-1}(u)=\tfrac12(\Tr(u)-u)$; it satisfies $\Tr\circ\psi=\Tr$. On $\mathrm{PG}(2,q)$ it acts
as the homology with axis $L$ ($\psi=-2\,\mathrm{id}$ on $\ker \Tr$) and centre $[1]$
($\psi(1)=1$); for $p=3$ one has $-2=1$ and $[1]\in L$, so $\psi$ is then an elation. One checks
that $W_{1/2}=\eta(-2)\cdot(W_0\circ\psi)$.
In conjugate coordinates $(y,y^q,y^{q^2})$, the cubic $\Nm(y)$ vanishes on the ``Singer
triangle'' of three conjugate lines. The map $\psi$ sends this triangle to its medial triangle,
so $W_{1/2}$ is $\eta$ of the quartic $(a+b+c)(-a+b+c)(a-b+c)(a+b-c)$ (Heron's product), while
$W_0$ is $\eta$ of $abc(a+b+c)$, up to constants.
\end{remark}

\begin{lemma}[Fourier form]\label{lem:jacobi}
Let $\chi$ be a character of $K^*$ that is trivial on $F^*$, i.e.\ a character of $G$. Then
\[
\chi(W_0)=\eta(-1)\sum_{y\in P}\eta_K(y)\chi(y),\qquad
\chi(W_{1/2})=\eta(2)\,J_P(\chi),
\]
where
\[
J_P(\chi):=\sum_{w\in K,\ \Tr(w)=1}\eta_K(w)\,\chi(1-w).
\]
Hence $W_{1/2}$ is a $\CW(m,q^2)$ if and only if $|J_P(\chi)|=q$ for every such $\chi$.
\end{lemma}
\begin{proof}
Use the representatives $y\in P$, for which $\Tr(y)=1$. Then
$W_{1/2}([y])=\eta(\Nm(\tfrac12-y))$. Substitute $w=1-2y$; this is a bijection of $P$ onto itself
since $\Tr(1)=3$. Then $\tfrac12-y=\tfrac12w$ and $\Nm(\tfrac12 w)=\Nm(w)/8$, so
$\eta(\Nm(\tfrac12w))=\eta(2)\eta_K(w)$. Moreover $\chi(y)=\chi(\tfrac12(1-w))=\chi(1-w)$,
because $\chi$ is trivial on $F^*$. The formula for $W_0$ is immediate. Finally apply
Lemma~\ref{lem:fourier}.
\end{proof}

The sum $J_P(\chi)$ is a Jacobi sum restricted to the hyperplane $\Tr=1$. The analogous
restricted Gauss sum $\sum_{y\in P}\eta_K(y)\chi(y)$ has absolute value $q$ by
Proposition~\ref{prop:classical}. We could not evaluate $J_P(\chi)$ in closed form. Numerically
(for $q=7$), $\chi(W_{1/2})/q$ is neither a root of unity times a Galois conjugate of
$\chi(W_0)$, nor a simple product or quotient of such values.

\begin{conjecture}[twin conjecture]\label{conj:twin}
For every odd prime power $q$, $W_{1/2}$ is a $\CW(q^2+q+1,q^2)$. Equivalently,
$|J_P(\chi)|=q$ for every character $\chi$ of $K^*$ trivial on $F^*$.
\end{conjecture}

\begin{theorem}\label{thm:twin}
Let $q$ be an odd prime power for which Conjecture~\ref{conj:twin} holds and $W_{1/2}\ne\pm
W_0$. Define $A\in\{0,\pm1\}^{2m}$ by
\[
A_x=\begin{cases}\tfrac12\bigl(W_0(x\bmod m)+W_{1/2}(x\bmod m)\bigr), & x\text{ even},\\[2pt]
\tfrac12\bigl(W_0(x\bmod m)-W_{1/2}(x\bmod m)\bigr), & x\text{ odd}.\end{cases}
\]
Then $A$ is a proper $\CW(2(q^2+q+1),q^2)$. The hypotheses hold, and hence such an $A$ exists,
for every odd prime power $q\le251$. For $q=11,13,17,19$ the resulting matrices
$\CW(266,11^2)$, $\CW(366,13^2)$, $\CW(614,17^2)$ and $\CW(762,19^2)$ are certified in Lean
(\S\ref{sec:lean}).
\end{theorem}
\begin{proof}
By Lemma~\ref{lem:support}, Proposition~\ref{prop:classical} and the hypothesis, $W_0$ and
$W_{1/2}$ are $\CW(m,q^2)$s with the same support, and $W_0$ is proper. Now apply
Corollary~\ref{cor:samesupport}. The verification for $q\le127$ is described in
Remark~\ref{rem:evidence}.
\end{proof}

\begin{remark}[computational evidence]\label{rem:evidence}
\leavevmode
\begin{enumerate}
\item For each of the 37 odd prime powers $q\le127$ we built $\F_{q^3}$ from a primitive
polynomial and computed $W_0$, $W_{1/2}$ and $A$. All periodic autocorrelations were computed
exactly. We used integer sequences and an FFT, checking that every value lay within $10^{-3}$ of
an integer before rounding; for $q\le19$ we also used direct $O(k^2)$ integer summation.
In every case $W_0$ and $W_{1/2}$ are $\CW(m,q^2)$s with the same support, $W_{1/2}\neq\pm W_0$,
and $A$ is a proper $\CW(2m,q^2)$. A second, independently written implementation
(\texttt{data/twin\_check\_gf.py}, which builds $\F_{q^3}$ as $\F_p[x]/(h)$ for every prime power
$q=p^k$ and uses its own primitive polynomials) confirms all of this for the 62 odd prime powers
$q\le251$ (\texttt{data/twin\_family\_check\_q\_le\_251.log}); for $q\le19$ its sequences are
multiplier-equivalent to the certified ones, and $W_{1/2}$ is never of the form $\pm W_0(tx+s)$.
\item For $q=3,5,7$, among all $\pm1$ sign patterns on $S$ that are constant on orbits of the
Frobenius multiplier $y\mapsto y^q$, the only ones giving a $\CW$ are $\pm W_0$ and
$\pm W_{1/2}$.
\item Among the $W_c$, $c\in F$, only $c\in\{0,\frac12\}$ give a $\CW$ for $q=5,7,11,13$.
\item For $q=5$, exactly $186$ of the $372000$ elements $g\in\mathrm{PGL}(3,5)$ make $W_0\circ g$
a $\CW$. Every $g\in N(C)\cup\psi N(C)$ does so, where $N(C)$ is the normaliser of the Singer
group $C$ ($|N(C)|=3m=93$). These two cosets are disjoint, because $\psi$ fixes all $q+1$ points of
$L$ while every non-identity element of $N(C)$ fixes at most $3$ points. Hence the $186$ elements
are exactly $N(C)\cup\psi N(C)$: up to the obvious symmetries, the twin is the only companion.
\item For the even prime power $q=8$, no Frobenius-invariant same-support companion of the
classical $\CW(73,64)$ exists (the construction needs $q$ odd).
\item A naive quartic-character analogue, intended to give $\CW(t(q^2+q+1),q^2)$ for
$t\mid q-1$ with $t>2$, fails for $t=4$ ($q=5,9,13$) and for $t=5,10$ ($q=11$).
\end{enumerate}
The certified rows for $q=11,13,17,19$ use primitive polynomials $x^3+10x^2+6x+9$,
$x^3+12x^2+9x+2$, $x^3+7x^2+8x+14$ and $x^3+2x^2+18x+4$ respectively (with $\alpha=x$). The file
\texttt{data/twin\_singer.py} rebuilds them from Definition~\ref{def:Wc} in pure Python and checks
equality with the certificates.
\end{remark}

\begin{remark}
For $q=3,5,7,9$, Theorem~\ref{thm:twin} produces proper $\CW(26,3^2)$, $\CW(62,5^2)$,
$\CW(114,7^2)$ and $\CW(182,9^2)$. In~\cite{G26} these orders appear as sporadic entries found by
earlier searches. We have not checked whether our matrices are equivalent to the tabulated ones.
In terms of~\cite[Theorem~6.6]{AGZ}, $2(q^2+q+1)=(q^3-1)/n$ with $n=(q-1)/2$. So the twin
construction re-establishes exactly this case of the withdrawn claim: unconditionally for the
certified values of $q$, and conjecturally in general. The other withdrawn orders, for example
$364=4\cdot91$ ($q=9$) and $732=4\cdot183$ ($q=13$), remain open.
\end{remark}

\section{A proper \texorpdfstring{$\CW(146,8^2)$}{CW(146,64)}}\label{sec:146}

Here $146=2\cdot73$ with $73=8^2+8+1$, but $q=8$ is even and \S\ref{sec:twin} does not apply.
Write a putative $\CW(146,64)$ as $A=U+zV$ as in Lemma~\ref{lem:doubling}. Then $B=U+V$ and
$B'=U-V$ satisfy $BB^{(-1)}=B'B'^{(-1)}=64$ and have entries in $\{0,\pm1,\pm2\}$. In other words,
they are integer weighing matrices $\ICW_2(73,64)$, with $B\equiv B'\pmod 2$ and
$|B\pm B'|\le2$ entrywise. Conversely, any such pair gives a $\CW(146,64)$ by
Lemma~\ref{lem:doubling}.

By McFarland's multiplier theorem in the form used in~\cite[Theorem~4.1]{AGZ}
($\gcd(73,64)=1$, $64=2^6$), $2$ is a multiplier of $B$ and of $B'$. We only used this to
restrict the search: we required $B$ to be fixed by $x\mapsto2x$ and $B'$ to be a translate of
a $2$-fixed element. The orbits of $x\mapsto2x$ on $\Z_{73}$ have sizes $1,9,9,\dots,9$.
Exhaustive enumeration of orbit-constant $B$ with entries in $\{0,\pm1,\pm2\}$, augmentation $8$
and $BB^{(-1)}=64$ gives $32$ solutions. Pairing them with all translates and signs and keeping
$\{0,\pm1\}$ lifts gives $112$ $\CW(146,64)$, of which $96$ are proper. In the certified
example, the positions where $B$ (respectively $B'$) equals $\pm2$ form a planar $(73,9,1)$
difference set, i.e.\ a line of $\mathrm{PG}(2,8)$. Unlike in \S\ref{sec:twin}, the supports of
$U$ and $V$ overlap. The first row is given in Appendix~\ref{app:data}.

We ran the same search for every repository-open entry $n=2m$ ($m$ odd) whose McFarland multiplier
group is nontrivial: $597$ entries, with at most $3\cdot10^5$ recursion nodes per entry. It found
only $n=146$. For $375$ entries it found no solution of this restricted form; the other $221$
entries (including $\CW(546,16^2)$, the $q=16$ analogue) reached the node limit. The negative
outcomes are \emph{not} nonexistence proofs.

\section{A proper \texorpdfstring{$\CW(504,12^2)$}{CW(504,144)}}\label{sec:504}

Let $G=\Z_{504}$, $H=21\Z_{504}\cong\Z_{24}$ and $K'=8\Z_{504}\cong\Z_{63}$. Then $H\cap K'=168\Z_{504}$
has order $3$ and $H+K'=G$. Take the proper $\CW(24,9)$ of the repository (due to Ang, Arasu, Ma and
Strassler) with $P=\{1,5,12,13,17,18\}$ and $N=\{0,6,9\}$, placed on $H$ via $i\mapsto21i$.
Take the proper $\CW(63,16)$ with $P=\{0,1,2,4,8,16,27,32,45,54\}$ and
$N=\{5,10,17,20,34,40\}$, placed on $K'$ via $j\mapsto8j$. The difference condition of
Lemma~\ref{lem:product} only needs checking inside $H\cap K'$; it holds there by direct
inspection (and by machine). So the product $\sum a_ib_j X^{21i+8j}$ is a proper $\CW(504,144)$.
Its divisors $d\ge144$ are $168$, $252$ and $504$. No proper $\CW(d,144)$ is known for any of
them (the repository lists $168$ as open and $252$ as nonexistent), so this is the first
$\CW(504,144)$ of any kind that we know of.
Iterating Lemma~\ref{lem:product} (and a generalized doubling) over all explicit matrices of
the repository produced no other pair $(n,k)$ missing from~\cite[Table~7]{G26}.

\section{Formal verification}\label{sec:lean}

The file \texttt{lean/Main.lean} (Lean~4.34.1, core library only, no Mathlib) contains the
following definitions (verbatim, except that line breaks in \texttt{MainStatement} are adjusted and
the two theorems are shown without their proofs). The statement uses the index form
\texttt{pafIdx}; the proofs use the fast form \texttt{paf}.
\begin{lstlisting}
def dot : List Int → List Int → Int
  | x :: xs, y :: ys => x * y + dot xs ys
  | _, _ => 0
def pafIdx (a : List Int) (d : Nat) : Int :=
  ((List.range a.length).map (fun i => a.getD i 0 * a.getD ((i + d) % a.length) 0)).sum
def paf (a : List Int) (d : Nat) : Int := dot a (a.rotateLeft d)
def IsCW (n k : Nat) (a : List Int) : Prop :=
  a.length = n ∧ (∀ x ∈ a, x = -1 ∨ x = 0 ∨ x = 1) ∧
  ∀ d, d < n → pafIdx a d = if d = 0 then (k : Int) else 0
def supp (a : List Int) : List Nat := (List.range a.length).filter (fun i => a.getD i 0 != 0)
def IsProper (a : List Int) : Prop :=
  ∀ e, e ≤ a.length → 1 < e → e ∣ a.length → ∃ i ∈ supp a, ∃ j ∈ supp a, i % e ≠ j % e
def MainStatement : Prop :=
    (∃ a, IsCW 146 64 a ∧ IsProper a) ∧ (∃ a, IsCW 266 121 a ∧ IsProper a) ∧
    (∃ a, IsCW 366 169 a ∧ IsProper a) ∧ (∃ a, IsCW 504 144 a ∧ IsProper a) ∧
    (∃ a, IsCW 614 289 a ∧ IsProper a) ∧ (∃ a, IsCW 762 361 a ∧ IsProper a)
theorem paf_eq_pafIdx (a : List Int) (d : Nat) (hd : d < a.length) :
    paf a d = pafIdx a d                      -- proved in general in the file
theorem new_proper_circulant_weighing_matrices : MainStatement   -- proved at the end
\end{lstlisting}
{\sloppy Each witness \texttt{cwX} (an explicit list of integers) comes with
\texttt{cwX\_isCW : IsCW n k cwX} and \texttt{cwX\_proper : IsProper cwX}. Both are proved by
\texttt{decide +kernel}. For the autocorrelation, the shift range is split into blocks of $64$
to bound kernel memory, the blocks are combined with \texttt{omega}, and the result is transported
to \texttt{pafIdx} by the general lemma \texttt{paf\_eq\_pafIdx}. That lemma is proved from
\texttt{dot\_eq\_sum} and \texttt{rotateLeft\_getD}. A second theorem,
\texttt{twin\_singer\_family\_members}, certifies the eight explicit rows of
Theorem~\ref{thm:twin} for $q=3,5,7,9,11,13,17,19$.\par}

\paragraph{Statement audit.}
\begin{itemize}
\item Since $n=$ \texttt{a.length}, all indices in \texttt{pafIdx} are $<n$, so \texttt{getD}
never returns its default. Thus \texttt{pafIdx a d} is literally
$\sum_{i<n}a_ia_{(i+d)\bmod n}$, and \texttt{IsCW} quantifies over all shifts $0\le d<n$, with
value $k$ at $d=0$ and $0$ otherwise. For the circulant matrix $W_{ij}=a_{(j-i)\bmod n}$ one has
$(WW^{\mathsf T})_{ij}=\PAF_a((i-j)\bmod n)$, so this is exactly $WW^{\mathsf T}=kI$. (This last
standard translation is not formalised.)
\item \texttt{IsProper}: the subgroups of $\Z_n$ are $e\Z_n$ with $e\mid n$, proper exactly when
$e>1$, and a set lies in a coset of $e\Z_n$ exactly when all its elements are congruent mod $e$.
The bound $e\le n$ is automatic for divisors.
\item The file contains sanity checks: $\CW(7,4)=-1+X+X^2+X^4$ satisfies
\texttt{IsCW 7 4} and \texttt{IsProper}; its lift to $\Z_{14}$ satisfies \texttt{IsCW 14 4} but
\emph{not} \texttt{IsProper}; and a corrupted sequence fails \texttt{IsCW}. So neither predicate
is vacuous.
\item The main statement has no hypotheses; it is a conjunction of closed existential claims with
explicit witnesses in the file.
\end{itemize}

\paragraph{Axioms and resources.} \texttt{lean/lean\_output.txt} reports, for the main theorem and
for every certificate theorem, dependence only on \texttt{propext}, \texttt{Classical.choice} and
\texttt{Quot.sound} (the \texttt{\_proper} theorems use only \texttt{propext}). There is no
\texttt{sorryAx} and no \texttt{native\_decide}. The full check took $491$\,s of wall time with a
maximum resident memory of $2.1$\,GB. The certificates were also checked independently by a
plain Python script (\texttt{data/verify\_certificates.py}).

\section{What is not verified; open questions}\label{sec:limits}

\begin{itemize}
\item \textbf{Conjecture~\ref{conj:twin} is open.} It is verified for odd prime powers $q\le251$
only, and the infinite family $\CW(2(q^2+q+1),q^2)$ is conditional on it. A proof would amount to
evaluating, or at least bounding exactly, the restricted Jacobi sums $J_P(\chi)$. Several
approaches failed for us: a direct Gauss-sum expansion leads to twisted character sums of cubic
polynomials; we found no collineation symmetry beyond Remark~\ref{rem:evidence}(4); and the
quartic analogue fails.
\item \textbf{Novelty} is judged against~\cite{GordonRepo,G26,AGZ,Tan26} and a literature search
on circulant weighing matrices and perfect ternary sequences. We could not consult
Krengel~\cite{Kr04}. The repository data file contains some stale entries (for example, some
orders covered by infinite families are still marked open), which is why we also compared with the
curated Table~7 of~\cite{G26}.
\item The finite-field constructions and the searches are not formalised. Lean certifies only the
final explicit sequences, which is all the main theorem needs.
\item Open: the even case ($q=16$ gives $\CW(546,16^2)$); the other withdrawn cases of
\cite[Theorem~6.6]{AGZ} ($364$ for $q=9$, $732$ for $q=13$); and the remaining open entries of
Strassler's table~\cite{AGZ,Tan26}, none of which is touched here.
\end{itemize}

\appendix
\section{First rows of the six new matrices}\label{app:data}

For each matrix we list $P=\{i:a_i=1\}$ and $N=\{i:a_i=-1\}$ (0-based positions in the first row
of the circulant; all other entries are $0$). The same data, together with the four
known-order family members, is in \texttt{data/certificates.json} and
\texttt{data/certificates\_PN.txt}.

\paragraph{$\mathrm{CW}(146,8^2)$ (Section~\ref{sec:146}).} Lean name \texttt{cw146}; $|P|=36$, $|N|=28$.\\
{\footnotesize\raggedright $P=\{$5, 7, 11, 15, 21, 22, 30, 34, 39, 42, 44, 47, 60, 66, 68, 83, 84, 87, 88, 90, 93, 94, 95, 101, 103, 106, 113, 115, 117, 118, 120, 126, 129, 132, 133, 136$\}$.\\
$N=\{$0, 2, 3, 4, 8, 16, 17, 19, 23, 32, 33, 45, 53, 59, 63, 64, 74, 79, 85, 97, 107, 110, 111, 119, 121, 128, 139, 141$\}$.\par}

\paragraph{$\mathrm{CW}(266,11^2)$ (Theorem~\ref{thm:twin}, $q=11$).} Lean name \texttt{cwFam266}; $|P|=66$, $|N|=55$.\\
{\footnotesize\raggedright $P=\{$4, 7, 8, 12, 14, 18, 24, 27, 28, 31, 32, 35, 38, 39, 40, 42, 44, 49, 50, 52, 58, 59, 75, 76, 77, 83, 86, 88, 91, 98, 102, 105, 106, 107, 113, 115, 117, 119, 120, 122, 132, 134, 144, 148, 152, 154, 156, 163, 170, 174, 179, 196, 197, 198, 201, 203, 213, 215, 218, 223, 237, 244, 245, 254, 256, 264$\}$.\\
$N=\{$0, 6, 9, 13, 16, 25, 34, 36, 45, 48, 57, 66, 69, 71, 72, 74, 79, 81, 93, 95, 99, 100, 103, 108, 118, 124, 125, 130, 138, 143, 162, 176, 180, 184, 186, 188, 189, 193, 194, 200, 206, 211, 217, 222, 225, 227, 229, 234, 243, 247, 249, 259, 260, 261, 262$\}$.\par}

\paragraph{$\mathrm{CW}(366,13^2)$ (Theorem~\ref{thm:twin}, $q=13$).} Lean name \texttt{cwFam366}; $|P|=91$, $|N|=78$.\\
{\footnotesize\raggedright $P=\{$2, 3, 12, 18, 20, 23, 25, 26, 31, 34, 35, 37, 39, 46, 48, 59, 60, 69, 76, 86, 88, 89, 95, 101, 105, 106, 107, 109, 112, 114, 115, 121, 129, 131, 137, 138, 141, 149, 156, 160, 165, 177, 179, 183, 190, 194, 197, 198, 199, 207, 212, 213, 215, 216, 227, 232, 233, 234, 239, 246, 249, 250, 256, 258, 260, 262, 264, 267, 268, 270, 274, 275, 280, 281, 293, 299, 309, 315, 317, 319, 322, 325, 326, 330, 338, 346, 353, 357, 358, 359, 365$\}$.\\
$N=\{$4, 5, 6, 9, 19, 21, 22, 36, 41, 42, 52, 53, 54, 57, 58, 65, 68, 70, 71, 78, 80, 93, 94, 96, 102, 111, 113, 117, 118, 120, 123, 124, 133, 135, 144, 146, 148, 150, 151, 152, 157, 164, 166, 167, 173, 178, 180, 181, 184, 191, 196, 211, 223, 228, 238, 247, 255, 265, 266, 273, 282, 283, 286, 287, 291, 302, 308, 310, 323, 328, 336, 337, 341, 342, 344, 345, 352, 355$\}$.\par}

\paragraph{$\mathrm{CW}(504,12^2)$ (Section~\ref{sec:504}).} Lean name \texttt{cw504}; $|P|=78$, $|N|=66$.\\
{\footnotesize\raggedright $P=\{$2, 4, 5, 21, 25, 29, 33, 37, 40, 53, 69, 80, 85, 90, 105, 108, 109, 113, 121, 129, 130, 136, 137, 149, 160, 166, 169, 180, 201, 206, 213, 229, 233, 234, 237, 252, 260, 262, 268, 269, 272, 273, 277, 281, 284, 285, 286, 289, 305, 306, 316, 320, 321, 325, 337, 349, 357, 361, 365, 373, 378, 380, 381, 386, 389, 394, 398, 401, 410, 421, 442, 446, 453, 461, 465, 468, 485, 489$\}$.\\
$N=\{$0, 8, 10, 13, 16, 20, 32, 34, 41, 45, 54, 61, 64, 68, 89, 101, 117, 125, 126, 128, 134, 142, 145, 146, 157, 158, 173, 181, 185, 189, 190, 194, 197, 205, 216, 221, 241, 253, 254, 256, 265, 292, 293, 313, 317, 332, 341, 342, 353, 360, 377, 382, 388, 397, 405, 409, 412, 418, 425, 432, 433, 437, 445, 458, 486, 493$\}$.\par}

\paragraph{$\mathrm{CW}(614,17^2)$ (Theorem~\ref{thm:twin}, $q=17$).} Lean name \texttt{cwFam614}; $|P|=153$, $|N|=136$.\\
{\footnotesize\raggedright $P=\{$2, 5, 6, 9, 10, 12, 13, 18, 27, 29, 34, 50, 55, 56, 60, 62, 73, 78, 84, 85, 88, 93, 98, 102, 104, 105, 107, 112, 118, 121, 122, 124, 129, 135, 141, 143, 145, 148, 153, 158, 159, 160, 163, 164, 167, 170, 173, 175, 176, 180, 189, 190, 200, 204, 207, 215, 217, 220, 221, 223, 224, 225, 226, 227, 230, 232, 236, 237, 240, 242, 247, 252, 258, 259, 260, 263, 264, 265, 266, 268, 273, 274, 286, 288, 290, 305, 306, 308, 315, 321, 324, 328, 330, 332, 333, 338, 339, 342, 343, 345, 351, 353, 359, 360, 371, 376, 378, 383, 396, 397, 398, 399, 406, 430, 434, 435, 438, 440, 441, 443, 444, 449, 453, 459, 475, 485, 493, 506, 515, 516, 519, 529, 536, 540, 545, 555, 556, 557, 564, 577, 578, 584, 585, 589, 591, 592, 594, 596, 598, 599, 600, 604, 609$\}$.\\
$N=\{$0, 4, 7, 19, 20, 22, 39, 41, 43, 47, 49, 57, 58, 61, 68, 75, 77, 79, 80, 81, 83, 94, 97, 101, 103, 108, 113, 115, 117, 119, 120, 125, 126, 132, 138, 140, 144, 147, 149, 150, 155, 156, 157, 169, 172, 179, 181, 183, 184, 185, 188, 195, 196, 198, 203, 210, 213, 214, 218, 219, 241, 245, 246, 254, 262, 267, 269, 275, 276, 283, 295, 296, 300, 310, 322, 323, 331, 335, 337, 340, 344, 352, 355, 358, 361, 370, 372, 374, 377, 381, 393, 394, 402, 403, 413, 417, 418, 421, 423, 437, 446, 458, 468, 478, 481, 484, 489, 498, 499, 500, 501, 504, 509, 512, 513, 518, 523, 535, 538, 541, 542, 550, 551, 558, 560, 562, 563, 568, 579, 586, 587, 588, 601, 606, 608, 611$\}$.\par}

\paragraph{$\mathrm{CW}(762,19^2)$ (Theorem~\ref{thm:twin}, $q=19$).} Lean name \texttt{cwFam762}; $|P|=190$, $|N|=171$.\\
{\footnotesize\raggedright $P=\{$0, 6, 7, 9, 11, 14, 29, 35, 43, 47, 53, 55, 56, 60, 66, 70, 78, 80, 83, 92, 94, 96, 97, 98, 99, 100, 103, 105, 114, 115, 124, 130, 131, 132, 133, 135, 138, 140, 142, 145, 147, 151, 153, 158, 161, 171, 174, 178, 184, 191, 193, 195, 197, 199, 201, 203, 204, 208, 209, 211, 212, 218, 219, 221, 222, 224, 227, 229, 241, 245, 248, 250, 251, 257, 258, 262, 266, 267, 271, 278, 279, 283, 286, 288, 291, 295, 300, 302, 303, 311, 319, 320, 324, 326, 330, 331, 332, 334, 336, 338, 344, 351, 357, 358, 360, 366, 367, 369, 371, 373, 374, 375, 376, 378, 382, 389, 396, 398, 400, 401, 404, 406, 408, 409, 412, 413, 417, 420, 422, 423, 433, 440, 443, 446, 448, 458, 460, 462, 469, 471, 482, 489, 492, 501, 503, 507, 529, 533, 536, 541, 551, 563, 567, 568, 573, 575, 577, 581, 583, 597, 607, 619, 621, 633, 642, 650, 657, 661, 665, 666, 675, 686, 687, 695, 704, 706, 710, 716, 720, 726, 727, 729, 733, 740, 742, 743, 744, 746, 758, 761$\}$.\\
$N=\{$3, 4, 12, 16, 18, 22, 24, 26, 45, 49, 50, 51, 54, 57, 58, 68, 71, 73, 76, 87, 93, 107, 109, 117, 119, 123, 129, 134, 136, 137, 139, 144, 154, 159, 163, 164, 165, 168, 169, 173, 175, 179, 181, 183, 185, 188, 198, 207, 213, 214, 228, 232, 233, 235, 237, 239, 242, 243, 249, 253, 255, 256, 260, 263, 264, 265, 270, 272, 273, 277, 282, 287, 292, 293, 297, 298, 304, 307, 309, 316, 317, 321, 322, 327, 328, 337, 340, 341, 342, 343, 353, 355, 364, 368, 379, 383, 386, 391, 402, 415, 418, 419, 421, 425, 429, 442, 444, 445, 450, 455, 456, 463, 466, 467, 470, 472, 476, 485, 487, 491, 494, 499, 506, 509, 522, 524, 527, 530, 531, 537, 538, 547, 548, 553, 557, 558, 561, 571, 586, 587, 591, 596, 598, 601, 611, 612, 615, 617, 625, 640, 649, 655, 656, 662, 670, 671, 677, 682, 689, 691, 693, 694, 696, 699, 714, 730, 731, 735, 737, 751, 753$\}$.\par}



\begin{thebibliography}{99}
\bibitem{AGZ} K.~T. Arasu, D.~M. Gordon, Y.~Zhang, \emph{New nonexistence results on circulant
weighing matrices}, Cryptogr.\ Commun.\ 13 (2021) 775--789; \href{https://arxiv.org/abs/1908.08447}{arXiv:1908.08447}.
\bibitem{AS98} K.~T. Arasu, J.~Seberry, \emph{On circulant weighing matrices}, Australas.\ J.\
Combin.\ 17 (1998) 21--37.
\bibitem{G26} D.~M. Gordon, \emph{Cyclic relative difference sets and circulant weighing
matrices}, J.\ Algebr.\ Comb.\ 63 (2026), Article 24; \href{https://arxiv.org/abs/2501.14924}{arXiv:2501.14924}.
\bibitem{GordonRepo} D.~M. Gordon, \emph{The La Jolla Circulant Weighing Matrix Repository},
data v1.3, commit 5b3d729 (2026-04-24),
\url{https://github.com/dmgordo/circulant-weighing-matrices};\\ table at
\url{https://dmgordon.org/circulant-weighing-matrices/} (site updated 2026-09-03; accessed
2026-10-08).
\bibitem{HJ83} T.~H{\o}holdt, J.~Justesen, \emph{Ternary sequences with perfect periodic
autocorrelation}, IEEE Trans.\ Inform.\ Theory 29 (1983) 597--600.
\bibitem{Kr04} E.~I. Krengel, \emph{Almost-perfect and odd-perfect ternary sequences}, in:
Sequences and Their Applications (SETA 2004), LNCS 3486, Springer, 2005, 197--207. (Not
consulted; only the abstract was seen.)
\bibitem{LS11} K.~H. Leung, B.~Schmidt, \emph{Finiteness of circulant weighing matrices of fixed
weight}, J.\ Combin.\ Theory Ser.\ A 118 (2011) 908--919.
\bibitem{Ma11} K.~H. Leung, S.~L. Ma, \emph{Proper circulant weighing matrices of weight $p^2$}, talk slides,
Institute for Mathematical Sciences, National University of Singapore (2011),
\url{https://imsarchives.nus.edu.sg/oldwww/Programs/011code/files/ma.pdf}.
\bibitem{SS13} B.~Schmidt, K.~W. Smith, \emph{Circulant weighing matrices whose order and weight
are products of powers of 2 and 3}, J.\ Combin.\ Theory Ser.\ A 120 (2013) 275--287.
\bibitem{Tan26} M.~M. Tan, \emph{Character and multiplier obstructions for circulant weighing
matrices}, \href{https://arxiv.org/abs/2608.18468}{arXiv:2608.18468} (2026).
\end{thebibliography}
\end{document}
